%% tex-Template for the 15th Viennese Conference on Optimal Control and Dynamic Games. 
%% ATTENTION: The Filename of this document MUST begin with the presenter's family name.
\documentclass[11pt,twoside,a4paper]{article}
%% Please try to avoid defining your own commands.
\usepackage{amsmath,amsfonts,amssymb}
\newcommand{\talktitle}[2]{\noindent {\em \hfill #2}\vspace{50pt}\newline\textbf{{#1}}\vspace{10pt}}
\newcommand{\speaker}[1]{\underline{#1}}
\newcommand{\coauthor}[3]{\noindent\mbox{#1 #2 (#3)}, }
\newcommand{\presenter}[3]{\noindent\mbox{\underline{#1 #2} (#3),} }
\newcommand{\affiliation}[2]{\noindent (#1)\,#2\;}
\renewcommand{\abstract}[1]{\vspace{20pt}#1}
\newcommand{\Skip}{\bigskip}
\textwidth=150mm

\begin{document}

%%Please enter the title and the authors of your talk. 
%\talktitle{Title}{Authors}
%%Please do not capitalize the title (except names) and  for the authors use 
%%initials like in the example here. Please distinguish the presenter by using
%\speaker{Author}
% Example title
\talktitle{Write the title of the abstract here (with uncapitalized words)}{M.~Mustermann, \speaker{A.N.~Other}, T.~Tester} 

%% Please enter the full name of the authors and the affiliation label.
%% For the presenting author please use the command
%   \presenter{First name middle names}{Last name}{affiliation label}
%% For the other authors please use
%   \coauthor{First name middle names}{Last name}{affiliation label}
%% such as follows
%% Example author 1
\coauthor{Max}{Mustermann}{1}
%% Example author 2
\presenter{Anne}{Other}{2}
%% Example author 3
\coauthor{Tommy}{Tester}{1}%

\Skip

%% Please enter the affiliations of the authors here 
% \affiliation{affiliation label}{Department, Institution, Town, Country}
%% Example affiliation 1
\affiliation{1}{VADOR, TU Wien, Vienna, Austria} 
%% Example affiliation 2
\affiliation{2}{BDA, University of Vienna, Vienna, Austria}



%% Text of the abstract:
% \abstract{Text}
%% Example abstract
\abstract{Please put the text of your abstract here. Keep the size of the 
abstract below half of a page. Use standard latex commands and try to avoid 
using your own defined commands. If you want to use enumerated formulas, please enumerate 
them manually with the use of the command {\tt $\backslash$eqno} or {\tt $\backslash$leqno}. 
You can see an example of how to use the command here:
\begin{displaymath}
a+b=c, \eqno (1)
\end{displaymath}
If you want to provide a bibliography please refer to the items manually using numbers, 
e.g. [1], and insert the bibliography at the end of the abstract using the environment 
{\tt description}, in the following style:  
\begin{description}
    \item[{[1]}] P. Grandits, R. M. Kovacevic, and V. M. Veliov. 
     Optimal control and the value of information for a stochastic epidemiological SIS-Model. 
   {\em Journal of Mathematical Analysis and Applications}, {\bf 476}(2):665-695, 2019.
     \item{[2]}
     R.M. Kovacevic and W. Semmler.
     Poverty trap and disaster insurance in a bi-level decision framework.   
     In {\em Dynamic Economic Problems with Regime Switches.}
     J. Haunschmied et al., Eds.
     Springer series {\em DMEEF}, vol. 25, 2020.
\end{description}
%% End of the text of the example abstract
}

\end{document}
